\section{improved quasi quark distribution through Wilson line renormalization}
Let us start by recalling the definition of the quasi quark distribution. For the unpolarized quark density, it is given as~\cite{Ji:2013dva}
\beq\label{qUnpolDef}
   \tilde q(x, \Lambda, p^z) = \int^\infty_{-\infty} \frac{dz}{4\pi} e^{izk^z}  \langle p|\overline{\psi}(0, 0_\perp, z)
   \gamma^z L(z,0)\psi(0) |p\rangle \ ,
\eeq
where the quark fields are separated along the spatial $z$-direction, $L(z,0)$ is the Wilson line gauge link inserted to ensure gauge invariance, and $\Lambda$ denotes the UV cutoff.

The one-loop correction to the above quark density has been computed both in the axial gauge~\cite{Xiong:2013bka} and in the Feynman gauge~\cite{Ji:2015jwa}. In Ref.~\cite{Ji:2015jwa}, the Wilson line $L(z,0)$ was separated into two Wilson lines along the $z$-axis $L(\infty,0)$ and $L(z,\infty)$, and associated respectively to the quark field to form gauge invariant combinations. Here we keep the gauge link as $L(z,0)$. The one-loop diagrams are then given in Fig.~\ref{1loopFeyn}. It is straightforward to check that these diagrams yield the same one-loop result as that in Ref.~\cite{Ji:2015jwa}. The linear divergence comes from the last diagram, which is a Wilson line self energy. It is known that such a linear divergence can be removed by a mass renormalization with the help of an auxiliary $z$-field formalism~\cite{Dotsenko:1979wb,Dorn:1986dt,Craigie:1980qs}, where the $z$-field is defined in a one-dimensional parameter space, and the non-local Wilson line can be interpreted as a two-point function of the $z$-field. In a sense, the auxiliary field $Z(z)$ can be understood as a Wilson line extending between $[z,\infty]$, i.e.
\beq
Z(z)=L(z,\infty),
\eeq
which satisfies an equation of motion 
\beq
\left[\partial_z-i g A_z(z)\right]Z(z)=0 ,
\eeq
analogous to a heavy quark field (in the heavy quark limit).   
The Wilson line 
\beq
L(z,0)=Z(z)Z^{\dagger}(0) 
\eeq
also renormalizes analogously to a heavy quark two point function 
\beq\label{WLrenorm}
L^{\text{ren}}(z,0)=\mathcal{Z}_Z^{-1}e^{-\delta m |z|}L(z,0) ,
\eeq
where $\delta m$ is analogous to the dimension one heavy quark mass counterterm and is linearly divergent, while the dimension zero wave function renormalization constant $\mathcal{Z}_Z$ is logarithmically divergent~\cite{Dotsenko:1979wb}. 


\begin{figure}[tbp]
\centering
\includegraphics[width=0.2\textwidth]{images/1loop1}
\hspace{2em}
\includegraphics[width=0.2\textwidth]{images/1loop2}
\hspace{2em}
\includegraphics[width=0.2\textwidth]{images/1loop3}
\hspace{2em}
\includegraphics[width=0.2\textwidth]{images/1loop4}
\caption{One-loop diagrams for quasi quark distribution in Feynman gauge, the conjugate diagrams are not shown.}
\label{1loopFeyn}
%\vspace*{-.5em}
\end{figure}

To illustrate the cancellation of linear divergence from the mass counterterm, let us determine $\delta m$ at one-loop level. Recall that the last diagram in Fig.~\ref{1loopFeyn} leads to the following linear divergence
\beq
\lim_{\epsilon\to 0}\int dk_z\frac{\alpha_s C_F\Lambda}{2\pi}\frac{[\delta(k_z-\bar x p_z)-\delta(\bar x p_z)]p_z}{k_z^2+\epsilon^2}
\eeq
with $\bar x=1-x$. On the other hand, plugging Eq.~(\ref{WLrenorm}) into Eq.~(\ref{qUnpolDef}), one finds that the mass counterterm yields the following contribution
\begin{align}
&-\int \frac{dz}{2\pi}p_z\, e^{i(x-1)p_z z}|z|\, \delta m=-\lim_{\epsilon\to 0}\int \frac{dz}{2\pi}p_z\, e^{-i \bar x p_z z}\frac{1-e^{-\epsilon |z|}}{\epsilon}\delta m\non\\
&=-\lim_{\epsilon\to 0}\int \frac{dz}{2\pi}p_z\, e^{-i \bar x p_z z}\int_0^\infty \frac{d\alpha}{\sqrt{\pi \alpha}}(1-e^{-\frac{z^2}{4\alpha}})e^{-\alpha \epsilon^2}\delta m\non\\
&=-\lim_{\epsilon\to 0}\int \frac{dz}{2\pi}p_z\, e^{-i \bar x p_z z}\int_0^\infty d\alpha \int d k_z \frac{e^{-\alpha (k_z^2+\epsilon^2)}(1-e^{i k_z z})}{\pi} \delta m\non\\
&=-\lim_{\epsilon\to 0}\int \frac{dz}{2\pi}p_z\, e^{-i \bar x p_z z}\int dk_z \frac{1}{\pi}\frac{(1-e^{i k_z z})}{k_z^2+\epsilon^2}\delta m\non\\
&=-\lim_{\epsilon\to 0}\int \frac{dk_z}{\pi}p_z \frac{\delta(\bar x p_z)-\delta(k_z-\bar x p_z)}{k_z^2+\epsilon^2}\delta m.
\end{align}
We therefore have
\beq
\delta m=-\frac{\alpha_s C_F}{2\pi}(\pi\Lambda)
\eeq
at one-loop. 

In coordinate space, by expanding the Wilson line to $\mathcal O(g^2)$, we have
\beq\label{WLcoordspace}
g^2 C_F\int_0^z dz_1\int_0^{z_1} dz_2 \frac{1}{4\pi^2}\frac{1}{(z_1-z_2)^2+a^2}=\frac{g^2 C_F}{4\pi^2}\big[\frac{z}{a}\tan^{-1}\big(\frac z a\big)-\frac 1 2\ln\big({1+\frac{z^2}{a^2}}\big)\big],
\eeq
where we have introduced a cutoff $a$ to regularize the short distance singularity in the coordinate space gluon propagator when $z_1\to z_2$. By Fourier transforming to momentum space, the above result leads to
\beq\label{eq8}
\frac{g^2 C_F}{8\pi^2}\big[\frac{1}{a p_z(1-x)^2}-\frac{1}{|1-x|}+C\delta(1-x)\big],
\eeq
where $C$ is a constant that can be determined by demanding the $x$-integration of the above expression vanish. We can therefore identify
\beq
\Lambda=\frac{1}{a}.
\eeq

The mass counterterm diagram of order $\alpha_s$ is shown in Fig.~\ref{1loopFeynmassCT}. The gauge independence of $\delta m$ can already be seen from that the one-loop correction in the axial gauge\B{~\cite{Xiong:2013bka}}  and Feynman gauge\B{~\cite{Ji:2015jwa}} contains identical linear divergences. By using the $R_\xi$ gauge, one easily finds that the extra piece $k^\mu k^\nu/k^2$ (where $k$ is the gluon momentum) in the numerator of gluon propagator has no impact on the linear divergence, therefore $\delta m$ remains the same in different gauges.


One can add more gluons and quarks to the one-loop diagrams in Fig.~\ref{1loopFeyn} to form higher loop diagrams. Power counting tells us that power divergence appears only in those diagrams which contain the Wilson line self energy as a subdiagram. As shown in the appendix, beyond one-loop the mass counterterm also removes all the power divergence from the Wilson line. Therefore, after including the mass counterterm contribution, there is at most logarithmic divergence in the quasi quark distribution. This indicates that, as far as the power divergence is concerned, the renormalization of a quark bilinear operator like the quasi quark distribution is the same as the renormalization of an open Wilson line. In the discussions above, we showed that the linear divergence is indeed cancelled by the mass counterterm at perturbative one-loop level. On the lattice, the mass counterterm $\delta m$ has to be determined nonperturbatively, which can be done e.g. as in Ref.~\cite{Musch:2010ka}.

\begin{figure}[tbp]
\centering
\includegraphics[width=0.2\textwidth]{images/1loopmassct}
\caption{One-loop mass counterterm diagram for quasi quark distribution in Feynman gauge.}
\label{1loopFeynmassCT}
\end{figure}

